
\documentclass{article}
\usepackage{colortbl}
\usepackage{makecell}
\usepackage{multirow}
\usepackage{supertabular}

\begin{document}

\newcounter{utterance}

\centering \large Interaction Transcript for game `cladder', experiment `full\_v1.5\_default', episode 1143 with Bentel rockers\_\_qwen3.5{-}4b\_\_Bentel\_iter.
\vspace{24pt}

{ \footnotesize  \setcounter{utterance}{1}
\setlength{\tabcolsep}{0pt}
\begin{supertabular}{c@{$\;$}|p{.15\linewidth}@{}p{.15\linewidth}p{.15\linewidth}p{.15\linewidth}p{.15\linewidth}p{.15\linewidth}}
    \# & \multicolumn{2}{c}{Player} && \multicolumn{2}{c}{Game Master} \\
    \hline

    \theutterance \stepcounter{utterance}  
    & & & \multicolumn{4}{p{0.6\linewidth}}{
        \cellcolor[rgb]{0.9,0.9,0.9}{
            \makecell[{{p{\linewidth}}}]{
                \texttt{\tiny{[P1$\langle$GM]}}
                \texttt{You are an expert in causal inference. The following question is not a typical commonsense query, but rather a meticulously designed question created by a professor specializing in causal inference, intended to assess the students' mastery of the course content.} \\
\texttt{Imagine a self{-}contained, hypothetical world with only the following conditions, and without any unmentioned factors or causal relationships: Kidney stone size has a direct effect on treatment and recovery. Treatment has a direct effect on recovery. Method 1: We look directly at how treatment correlates with recovery in general. Method 2: We look at this correlation case by case according to kidney stone size. To understand how treatment affects recovery, is it more correct to use the Method 1 than Method 2?} \\
\texttt{Start your answer with 'Yes' or 'No', followed by additional reasoning or evidence to support your explanation.} \\
            }
        }
    }
     \\ \\

    \theutterance \stepcounter{utterance}  
    & \multicolumn{4}{p{0.6\linewidth}}{
        \cellcolor[rgb]{0.561,0.737,0.561}{
            \makecell[{{p{\linewidth}}}]{
                \texttt{\tiny{[P1$\rangle$GM]}}
                \texttt{No} \\
\\ 
\texttt{**Reasoning:**} \\
\\ 
\texttt{In this hypothetical world, **Kidney stone size** acts as a **confounding variable** (or a common cause) for both the **Treatment** and the **Recovery**.} \\
\texttt{1.  **The Structure**: We have a causal path where Kidney Stone Size \$\textbackslash{}rightarrow\$ Treatment and Kidney Stone Size \$\textbackslash{}rightarrow\$ Recovery. Additionally, there is a direct causal path Treatment \$\textbackslash{}rightarrow\$ Recovery.} \\
\texttt{2.  **The Problem with Method 1**: When we look at the correlation between Treatment and Recovery in general (without conditioning on stone size), we are observing the **backdoor path** (Recovery \$\textbackslash{}leftarrow\$ Treatment \$\textbackslash{}leftarrow\$ Kidney Stone Size \$\textbackslash{}rightarrow\$ Recovery). This path induces a spurious correlation between Treatment and Recovery because patients with larger stones are more likely to receive different treatments and have different recovery rates independently of the treatment itself. Therefore, the raw correlation in Method 1 is biased and does not reflect the true causal effect of treatment on recovery.} \\
\texttt{3.  **The Solution with Method 2**: By analyzing the correlation case by case according to kidney stone size, we are effectively **conditioning on the confounder** (stratification). This blocks the backdoor path. Within each stratum (e.g., "small stones" or "large stones"), the distribution of stone sizes is constant, so the spurious correlation induced by stone size is eliminated. The remaining correlation between Treatment and Recovery within these strata represents the **local causal effect** for that specific group.} \\
\texttt{4.  **Conclusion**: To isolate the true causal effect of treatment on recovery in the presence of a confounder, one must adjust for the confounder. Method 2 achieves this adjustment, whereas Method 1 does not. Therefore, Method 2 is the correct approach to understand how treatment affects recovery.} \\
            }
        }
    }
    & & \\ \\

    \theutterance \stepcounter{utterance}  
    & & & \multicolumn{2}{p{0.3\linewidth}}{
        \cellcolor[rgb]{0.9,0.9,0.9}{
            \makecell[{{p{\linewidth}}}]{
                \texttt{\tiny{[GM$|$GM]}}
                \texttt{No} \\
            }
        }
    }
    & & \\ \\

    \theutterance \stepcounter{utterance}  
    & & & \multicolumn{2}{p{0.3\linewidth}}{
        \cellcolor[rgb]{0.9,0.9,0.9}{
            \makecell[{{p{\linewidth}}}]{
                \texttt{\tiny{[GM$|$GM]}}
                \texttt{no} \\
            }
        }
    }
    & & \\ \\

    \theutterance \stepcounter{utterance}  
    & & & \multicolumn{2}{p{0.3\linewidth}}{
        \cellcolor[rgb]{0.9,0.9,0.9}{
            \makecell[{{p{\linewidth}}}]{
                \texttt{\tiny{[GM$|$GM]}}
                \texttt{game\_result = WIN} \\
            }
        }
    }
    & & \\ \\

\end{supertabular}
}

\end{document}
